\documentclass[a4paper,12pt]{article}
\usepackage[utf8]{inputenc}
\usepackage[english]{babel}
\usepackage{listings}
\usepackage{xcolor}
\usepackage{float}
\usepackage{placeins}

\lstset{
    language=C,
    backgroundcolor=\color{lightgray!20},
    frame=single,
    breaklines=true,
    numbers=left,
    numberstyle=\tiny,
    stepnumber=1,
    numbersep=5pt,
    captionpos=b,
    basicstyle=\ttfamily\small,
    keywordstyle=\color{blue}\bfseries,
    commentstyle=\color{green!60!black},
    stringstyle=\color{red},
    showstringspaces=false
}

\title{LabReSiD25 - Hands-On 1}
\author{Davide Ferrara - 518629}
\date{}

\begin{document}

\maketitle

\section{Esercizio 1: Comunicazione UDP}

\subsection{Obbiettivo}
Usando i Linux Socket scrivere un server (\texttt{server\_udp.c}) ed un client
(\texttt{client\_udp.c}) che realizzino una comunicazione bidirezionale 
basata su UDP.
\newline
UDP (User Datagram Protocol) è un protocollo non orientato alla
connessione del modello TCP/IP. A differenza di TCP non richiede
\texttt{connect()}/\texttt{listen()}/\texttt{accept()} ma utilizza direttamente
\texttt{recvfrom()} e \texttt{sendto()}.

\subsection{Server UDP}

Includo le librerie necessarie per la comunicazione di rete.
\begin{lstlisting}[caption={server\_udp.c - Include}]
#include <arpa/inet.h>   // inet_ntoa, htons, ntohs
#include <stddef.h>      // size_t
#include <stdio.h>       // printf, perror
#include <stdlib.h>      // exit
#include <string.h>      // memset, strlen
#include <sys/socket.h>  // socket, bind, recvfrom, sendto
#include <unistd.h>       // close
\end{lstlisting}
\FloatBarrier

Definisco le costanti per la porta e la dimensione del buffer.
\begin{lstlisting}[caption={server\_udp.c - Costanti}]
#define PORT 8080
#define BUF_SIZE 1024
\end{lstlisting}
\FloatBarrier

Creo la struttura \texttt{sockaddr\_in} per IPv4 contenente indirizzo e porta.
\begin{itemize}
    \item \texttt{AF\_INET}: famiglia di indirizzi IPv4
    \item \texttt{htons()}: converte la porta in formato di rete (big-endian)
    \item \texttt{INADDR\_ANY = 0.0.0.0}: il server ascolta su TUTTE le interfacce di rete
\end{itemize}
\begin{lstlisting}[caption={server\_udp.c - sockaddr\_in}]
struct sockaddr_in addr_in;
memset(&addr_in, 0, sizeof(addr_in));
addr_in.sin_family = AF_INET;
addr_in.sin_port = htons(PORT);
addr_in.sin_addr.s_addr = INADDR_ANY;
\end{lstlisting}
\FloatBarrier

Creo il socket di tipo \texttt{SOCK\_DGRAM} (datagram) per UDP. La \texttt{socket()} restituisce un file descriptor. Controllo con \texttt{-1} per gestire gli errori.
\begin{lstlisting}[caption={server\_udp.c - socket()}]
int server_fd = socket(AF_INET, SOCK_DGRAM, 0);
if (server_fd == -1) {
    perror("Creazione socket fallita!");
    exit(1);
}
\end{lstlisting}
\FloatBarrier

Il \texttt{bind()} associa il socket all'indirizzo e porta specificati nella struttura. Il server rimane in ascolto su quella porta.
\begin{lstlisting}[caption={server\_udp.c - bind()}]
if (bind(server_fd, (struct sockaddr *)&addr_in, sizeof(addr_in)) == -1) {
    perror("Errore nel bind");
    close(server_fd);
    exit(1);
}
printf("Server UDP in ascolto sulla porta %d...\n", PORT);
\end{lstlisting}
\FloatBarrier

Alloco il buffer per ricevere i dati e la struttura \texttt{sockaddr\_in} per memorizzare le informazioni sul client.
\begin{lstlisting}[caption={server\_udp.c - Buffer}]
char buf[BUF_SIZE];
struct sockaddr_in client_addr;
socklen_t client_len = sizeof(client_addr);
\end{lstlisting}
\FloatBarrier

Ciclo infinito: \texttt{recvfrom()} è bloccante e riceve datagrammi UDP,
restituisce anche l'indirizzo del client e dopo aver stampato il messaggio,
rispondo con \texttt{sendto()}.

\begin{lstlisting}[caption={server\_udp.c - recvfrom/sendto}]
while (1) {
    ssize_t n = recvfrom(server_fd, buf, sizeof(buf) - 1, 0,
                         (struct sockaddr *)&client_addr, &client_len);
    if (n > 0) {
        buf[n] = '\0';
        printf("Ricevuto da %s:%d: %s\n", 
               inet_ntoa(client_addr.sin_addr),
               ntohs(client_addr.sin_port), buf);
        char *msg = "pong\n";
        sendto(server_fd, msg, strlen(msg), 0,
               (struct sockaddr *)&client_addr, client_len);
    }
}
\end{lstlisting}
\FloatBarrier

\subsection{Client UDP}

Struttura dati analoga al server, ma per il server remoto,
uso \texttt{127.0.0.1} (localhost) come indirizzo del server.
\begin{lstlisting}[caption={client\_udp.c - sockaddr\_in}]
struct sockaddr_in server_addr;
memset(&server_addr, 0, sizeof(server_addr));
server_addr.sin_family = AF_INET;
server_addr.sin_port = htons(PORT);
server_addr.sin_addr.s_addr = inet_addr("127.0.0.1");
\end{lstlisting}
\FloatBarrier

\newpage
Invio il messaggio "ping"" al server con \texttt{sendto()},
ricevo la risposta con \texttt{recvfrom()}.
Il \texttt{NULL} nei parametri finali indica che non mi serve l'indirizzo
del server (lo conosco già).
\newline
\texttt{sleep(1)} invece pausa il ciclo di 1 secondo.
\begin{lstlisting}[caption={client\_udp.c - Invio/Ricezione}]
char buf[BUF_SIZE];
while (1) {
    sendto(client_fd, "ping\n", 5, 0,
           (struct sockaddr *)&server_addr, sizeof(server_addr));
    printf("Inviato a %s:%d: ping\n", ADDR, PORT);
    
    ssize_t n = recvfrom(client_fd, buf, sizeof(buf) - 1, 0, NULL, NULL);
    if (n > 0) {
        buf[n] = '\0';
        printf("Ricevuto da %s:%d: %s\n", ADDR, PORT, buf);
    }
    sleep(1);
}
\end{lstlisting}
\FloatBarrier

Chiudo il file descriptor.
\begin{lstlisting}[caption={client\_udp.c - Chiusura}]
close(client_fd);
return 0;
\end{lstlisting}
\FloatBarrier

\subsection{Differenze Client/Server}
Il server effettua il \texttt{bind()} sulla porta per rimanere in ascolto. Il client invia direttamente al server senza bind. La restante logica è praticamente uguale.

\subsection{Output di esempio}
\begin{lstlisting}[caption={Output server}]
Server UDP in ascolto sulla porta 8080...
Ricevuto da 127.0.0.1:37698: ping
Ricevuto da 127.0.0.1:37698: ping
Ricevuto da 127.0.0.1:37698: ping
\end{lstlisting}
\FloatBarrier

\begin{lstlisting}[caption={Output client}]
Avvio del Client...
Inviato a 127.0.0.1:8080: ping
Ricevuto da 127.0.0.1:8080: pong
\end{lstlisting}
\FloatBarrier

\section{Compilazione}
Per compilare entrambi i programmi usiamo \texttt{make}.
\begin{lstlisting}[caption={Compilazione}]
make
\end{lstlisting}
\FloatBarrier
Per rimuovere i file compilati usiamo \texttt{make clean}.
\begin{lstlisting}[caption={Clean}]
make clean
\end{lstlisting}
\FloatBarrier

\section{Esercizio 2: File Descriptor}

\subsection{Obbiettivo}
Listare, attraverso l'opportuno tool da riga di comando, il file descriptor
relativo al programma precedente ed analizzarne l'output.

\subsection{Tool: lsof}
Il comando \texttt{lsof} (List Open Files) permette di visualizzare i file aperti da un processo.
\begin{lstlisting}[caption={Comando lsof}]
bash$ lsof -i :8080
\end{lstlisting}
\FloatBarrier

\newpage
\subsection{Output}
\begin{lstlisting}[caption={Output lsof}]
COMMAND     PID  USER FD   TYPE DEVICE SIZE/OFF NODE NAME
server_udp 70928 dff    3u  IPv4 334618  0t0    UDP *:http-alt
\end{lstlisting}
\FloatBarrier

\subsection{Analisi}
\begin{itemize}
    \item \textbf{COMMAND}: nome del comando/processo (\texttt{server\_udp})
    \item \textbf{PID}: identificativo univoco del processo (70928)
    \item \textbf{USER}: utente che ha lanciato il processo (dff)
    \item \textbf{FD 3u}: file descriptor 3 in modalita' read/write

    \textbf{Contesto sui file descriptor:}
    Ogni processo in Linux ha i suoi file descriptor. I primi 3 sono riservati:
    \begin{itemize}
        \item fd 0: standard input (stdin)
        \item fd 1: standard output (stdout)
        \item fd 2: standard error (stderr)
    \end{itemize}
    I socket di rete vengono assegnati ai fd successivi. Il nostro server ha il socket su fd 3. La lettera \texttt{u} indica che il socket UDP supporta sia invio che ricezione.

    \item \textbf{TYPE IPv4}: famiglia di indirizzi IPv4
    \item \textbf{DEVICE}: numero del nodo di rete assegnato dal kernel
    \item \textbf{SIZE/OFF}: dimensione del buffer o offset (0t0 = nessun offset)
    \item \textbf{NODE}: numero di porta assegnato dal kernel (334618)
    \item \textbf{NAME}: indirizzo di binding
\end{itemize}

\section{Conclusioni}
In questa hands-on ho implementato un sistema client-server UDP funzionante,
comprendendo le differenze tra socket UDP e TCP. Ho inoltre appreso come
monitorare i file descriptor dei processi di rete utilizzando \texttt{lsof}.

\end{document}
